\documentclass[11pt]{article}
\usepackage[utf8]{inputenc}
\usepackage[margin=1in]{geometry}
\usepackage{amsmath}
\usepackage{graphicx}
\usepackage{booktabs}
\usepackage{hyperref}
\usepackage{natbib}
\usepackage{lineno}
\usepackage{xcolor}
\usepackage{multirow}
\usepackage{float}

\linenumbers

\title{Edge-of-Chaos Criticality Mediates Cross-Frequency Cortical--Thalamic Information Transfer During Consciousness}

\author{
Integrative Neurophysiology Laboratory\\
Department of Systems Neuroscience\\
}

\date{}

\begin{document}

\maketitle

\begin{abstract}
Consciousness requires preserved bidirectional communication between cortex and thalamus, yet the spectral mechanisms of this communication and the dynamical principles governing their modulation across brain states remain poorly understood. Here, using spectrally resolved directed transfer entropy in simultaneous thalamic--cortical recordings from humans, rats, and mice, we identified a highly conserved cross-frequency communication channel: information encoded in low-frequency oscillations (1.5--13~Hz) at the sender is decoded by high gamma activity (52--104~Hz) at the receiver. This channel was significantly reduced during propofol anesthesia (humans: $W = 3$, $p = 0.016$; rats: $W = 0$, $p = 0.008$; one-tailed Wilcoxon signed-rank) and during generalized spike-and-wave seizures ($W = 1$, $p = 0.012$), while the psychedelic 5-MeO-DMT selectively enhanced cortex-to-thalamus transfer ($W = 15$, $p = 0.031$). Non-spectrally-resolved transfer entropy and phase-amplitude coupling showed no consistent relationship with consciousness across species or conditions. A Bayesian-genetic optimized mean-field model of the basal ganglia--thalamocortical system demonstrated that proximity to edge-of-chaos criticality of slow thalamocortical dynamics governs the strength of cross-frequency information transfer, with anesthesia moving dynamics away from the critical point (Lyapunov exponent shift: $\Delta\lambda = -0.18 \pm 0.04$) and psychedelics tuning dynamics closer ($\Delta\lambda = +0.09 \pm 0.03$). These findings reveal a mechanistic link between criticality, cross-frequency neural communication, and consciousness.
\end{abstract}

\section{Introduction}

Consciousness is accompanied by rich, bidirectional information exchange between the thalamus and the cortex \citep{tononi2016integrated, mashour2020conscious}. The thalamus relays sensory information to cortex and receives extensive feedback projections, forming recurrent loops essential for sustaining conscious experience. Disruption of this thalamocortical communication---whether by general anesthesia \citep{alkire2008consciousness}, seizures, or structural lesions---reliably produces loss of consciousness, underscoring the centrality of these circuits to awareness.

Despite the established importance of thalamocortical connectivity, the spectral organization of directed information flow between these structures has received limited attention. Neural oscillations at different frequencies serve distinct computational roles \citep{canolty2010importance}, and cross-frequency interactions enable the coordination of local and global processing \citep{palva2005neural}. Whether specific frequency-pair combinations carry preferential information between cortex and thalamus, and whether such channels are modulated by brain state, has not been systematically investigated across species.

The critical brain hypothesis provides a compelling framework for understanding state-dependent neural dynamics \citep{beggs2003neuronal, shew2013functional}. Systems operating near the edge of chaos---the boundary between stable, ordered dynamics and unstable, chaotic dynamics---exhibit maximal information processing capacity, sensitivity, and dynamic range. The Lyapunov exponent quantifies proximity to this critical point: negative values indicate stable dynamics, positive values indicate chaos, and the transition between them marks the edge of chaos. Whether this dynamical regime governs the efficiency of thalamocortical communication has not been tested.

Here we address these gaps through a multi-species investigation combining spectrally resolved directed transfer entropy \citep{schreiber2000measuring} with mean-field computational modeling \citep{vanalbada2009mean}. We recorded simultaneous thalamic and cortical electrophysiology from humans, rats, and mice during waking, anesthesia, seizures, and psychedelic states, and we characterized how edge-of-chaos criticality shapes cross-frequency information transfer.

\section{Results}

\subsection{Identification of a Conserved Cross-Frequency Communication Channel}

An initial exploratory sweep of spectrally resolved directed transfer entropy across all frequency-pair combinations revealed a prominent motif of low-to-high frequency communication in all four species and recording configurations during the waking state. Z-scored information transfer matrices showed elevated transfer from delta/theta/alpha band (1.625--13~Hz) oscillations at the sender to high gamma (52--104~Hz) oscillations at the receiver, in both the cortex-to-thalamus and thalamus-to-cortex directions.

\begin{table}[H]
\centering
\caption{Surrogate-tested significance of cross-frequency transfer entropy (low 1.5--13~Hz to high 52--104~Hz) during waking across species. Surrogate testing used frequency-specific randomization of dynamics; significance assessed per 10-second window. Fraction significant indicates proportion of windows reaching $p < 0.05$ after surrogate correction.}
\label{tab:surrogate}
\begin{tabular}{llcccc}
\toprule
Species & Direction & $n$ (subjects) & Fraction Sig. & Mean TE (bits) & SE \\
\midrule
Human (ET) & Ctx $\to$ Thal & 8 & 0.82 & 0.047 & 0.008 \\
Human (ET) & Thal $\to$ Ctx & 8 & 0.78 & 0.041 & 0.007 \\
Long-Evans rat & Ctx $\to$ Thal & 6 & 0.89 & 0.053 & 0.011 \\
Long-Evans rat & Thal $\to$ Ctx & 6 & 0.85 & 0.048 & 0.009 \\
C57/BL6 mouse & Ctx $\to$ Thal & 5 & 0.91 & 0.062 & 0.014 \\
C57/BL6 mouse & Thal $\to$ Ctx & 5 & 0.87 & 0.055 & 0.012 \\
GAERS rat & Ctx $\to$ Thal & 7 & 0.76 & 0.039 & 0.010 \\
GAERS rat & Thal $\to$ Ctx & 7 & 0.72 & 0.035 & 0.009 \\
\bottomrule
\end{tabular}
\end{table}

Rigorous surrogate testing confirmed that this cross-frequency transfer was statistically significant in all species and both directions (Table~\ref{tab:surrogate}). The fraction of 10-second windows reaching significance ranged from 0.72 to 0.91, with mean transfer entropy values of 0.035--0.062 bits. This consistency across species with different thalamic nuclei and cortical areas (Table~\ref{tab:recording}) supports a highly conserved spectral organization of thalamocortical communication.

\begin{table}[H]
\centering
\caption{Electrophysiology recording configurations across species. All recordings were simultaneous thalamic--cortical pairs. Connection type indicates known anatomical pathway between recorded structures.}
\label{tab:recording}
\begin{tabular}{llllc}
\toprule
Species & Thalamic Nucleus & Cortical Region & Connection & $n$ \\
\midrule
Human (ET) & Vim & Ipsilateral SMC & Direct reciprocal & 8 \\
Long-Evans rat & VP & Ipsilateral S1 & Direct reciprocal & 6 \\
C57/BL6 mouse & MD & mPFC & Direct reciprocal & 5 \\
GAERS rat & VP & Contralateral S1 & Indirect & 7 \\
\bottomrule
\end{tabular}
\end{table}

\subsection{Unconscious States Reduce Cross-Frequency Transfer}

\begin{table}[H]
\centering
\caption{Effect of brain state on cortex-to-thalamus cross-frequency transfer entropy. One-tailed Wilcoxon signed-rank tests compare waking vs.\ altered state within each species. Effect size $r = Z / \sqrt{n}$. $\Delta$TE represents mean change from waking baseline. $^{*}p<0.05$, $^{**}p<0.01$.}
\label{tab:ctx_thal}
\begin{tabular}{llccccc}
\toprule
Comparison & Species & $n$ & $\Delta$TE (bits) & $W$ & $p$ & $r$ \\
\midrule
Wake vs.\ Propofol & Human ET & 8 & $-0.019$ & 3 & $0.016^{*}$ & 0.73 \\
Wake vs.\ Propofol & LE rat & 6 & $-0.028$ & 0 & $0.008^{**}$ & 0.88 \\
Wake vs.\ Seizure & GAERS rat & 7 & $-0.022$ & 1 & $0.012^{*}$ & 0.81 \\
Wake vs.\ 5-MeO-DMT & C57 mouse & 5 & $+0.024$ & 15 & $0.031^{*}$ & 0.78 \\
\bottomrule
\end{tabular}
\end{table}

Propofol anesthesia significantly reduced cortex-to-thalamus cross-frequency transfer in both humans ($W = 3$, $p = 0.016$, $r = 0.73$) and rats ($W = 0$, $p = 0.008$, $r = 0.88$; Table~\ref{tab:ctx_thal}). Generalized spike-and-wave seizures in GAERS rats similarly reduced this transfer ($W = 1$, $p = 0.012$, $r = 0.81$). Conversely, the psychedelic 5-MeO-DMT significantly enhanced cortex-to-thalamus transfer in mice ($W = 15$, $p = 0.031$, $r = 0.78$), consistent with the expanded consciousness hypothesis of psychedelic states \citep{carhart2016neural}.

\begin{table}[H]
\centering
\caption{Effect of brain state on thalamus-to-cortex cross-frequency transfer entropy. Same statistical framework as Table~\ref{tab:ctx_thal}.}
\label{tab:thal_ctx}
\begin{tabular}{llccccc}
\toprule
Comparison & Species & $n$ & $\Delta$TE (bits) & $W$ & $p$ & $r$ \\
\midrule
Wake vs.\ Propofol & Human ET & 8 & $-0.017$ & 4 & $0.023^{*}$ & 0.68 \\
Wake vs.\ Propofol & LE rat & 6 & $-0.025$ & 1 & $0.016^{*}$ & 0.82 \\
Wake vs.\ Seizure & GAERS rat & 7 & $-0.018$ & 2 & $0.018^{*}$ & 0.76 \\
Wake vs.\ 5-MeO-DMT & C57 mouse & 5 & $+0.011$ & 12 & $0.156$ & 0.44 \\
\bottomrule
\end{tabular}
\end{table}

Thalamus-to-cortex transfer was also significantly reduced during unconscious states (Table~\ref{tab:thal_ctx}). Notably, 5-MeO-DMT did not significantly enhance thalamus-to-cortex transfer ($W = 12$, $p = 0.156$), suggesting a directional asymmetry in psychedelic modulation of thalamocortical communication, possibly due to limited statistical power with $n = 5$ animals.

\subsection{Broadband Transfer Entropy and Phase-Amplitude Coupling Show No Consistent Pattern}

\begin{table}[H]
\centering
\caption{Control analyses: non-spectrally-resolved transfer entropy and phase-amplitude coupling (modulation index) for cortex-to-thalamus direction. Neither measure shows a consistent relationship with consciousness across species and conditions. Wilcoxon signed-rank tests; consistency coded as whether direction of change aligns with consciousness level.}
\label{tab:controls}
\begin{tabular}{llcccccc}
\toprule
 & & \multicolumn{3}{c}{Broadband TE} & \multicolumn{3}{c}{PAC (MI)} \\
\cmidrule(lr){3-5} \cmidrule(lr){6-8}
Comparison & Species & $\Delta$ & $p$ & Consistent? & $\Delta$ & $p$ & Consistent? \\
\midrule
Wake vs.\ Propofol & Human & $+0.003$ & 0.461 & No & $-0.008$ & 0.195 & No \\
Wake vs.\ Propofol & LE rat & $-0.011$ & 0.078 & No & $+0.012$ & 0.109 & No \\
Wake vs.\ Seizure & GAERS & $+0.007$ & 0.312 & No & $-0.005$ & 0.406 & No \\
Wake vs.\ 5-MeO & Mouse & $-0.002$ & 0.812 & No & $+0.003$ & 0.688 & No \\
\bottomrule
\end{tabular}
\end{table}

Neither broadband (non-spectrally-resolved) transfer entropy nor phase-amplitude coupling showed consistent relationships with consciousness level across species (Table~\ref{tab:controls}). Broadband TE estimated using the Kraskov nearest-neighbor method \citep{kraskov2004estimating} via the Java Information Dynamics Toolkit \citep{lizier2014jidt} produced variable and non-significant results. The modulation index for cross-frequency PAC \citep{tort2010measuring} likewise showed no reliable pattern. These null results underscore the critical importance of spectral resolution in characterizing consciousness-related communication.

\subsection{Power Spectral Density Changes Do Not Explain Transfer Entropy Results}

\begin{table}[H]
\centering
\caption{Power spectral density changes across brain states (Welch's method, median across subjects). Both propofol and 5-MeO-DMT increase slow power and decrease high-frequency power, yet have opposing effects on cross-frequency TE, ruling out simple spectral power explanations. Values in dB relative to waking baseline.}
\label{tab:psd}
\begin{tabular}{llcccc}
\toprule
Manipulation & Region & $\Delta$ Slow ($\leq$4~Hz) & $\Delta$ High ($>$80~Hz) & $\Delta$ Cross-Freq TE & Dissociation? \\
\midrule
Propofol & Cortex & $+4.2$ dB & $-3.8$ dB & $-0.019$ bits & --- \\
Propofol & Thalamus & $+3.7$ dB & $-4.1$ dB & $-0.019$ bits & --- \\
5-MeO-DMT & Cortex & $+2.9$ dB & $-2.1$ dB & $+0.024$ bits & Yes \\
5-MeO-DMT & Thalamus & $+3.1$ dB & $-2.4$ dB & $+0.024$ bits & Yes \\
Seizure & Cortex & $+6.8$ dB & $-1.2$ dB & $-0.022$ bits & --- \\
Seizure & Thalamus & $+5.4$ dB & $-1.8$ dB & $-0.022$ bits & --- \\
\bottomrule
\end{tabular}
\end{table}

Both propofol and 5-MeO-DMT increased slow-frequency power and decreased high-frequency power (Table~\ref{tab:psd}), yet they had opposing effects on cross-frequency transfer entropy ($-0.019$ vs.\ $+0.024$ bits). This dissociation demonstrates that spectral power changes alone cannot account for the observed modulation of directed information transfer, pointing instead to a dynamical mechanism governing the efficiency of cross-frequency communication.

\subsection{Computational Model Links Criticality to Cross-Frequency Transfer}

\begin{table}[H]
\centering
\caption{Edge-of-chaos criticality and cross-frequency transfer in the mean-field model. The Lyapunov exponent ($\lambda$) quantifies proximity to the edge of chaos ($\lambda = 0$). Model parameters optimized via Bayesian-genetic methods to match empirical waking electrodynamics. $\Delta\lambda$ computed relative to optimized waking state. Mean $\pm$ SD across 100 parameter perturbation samples.}
\label{tab:model}
\begin{tabular}{lcccc}
\toprule
Brain State & $\lambda$ & $\Delta\lambda$ & Model TE (bits) & $\Delta$TE \\
\midrule
Waking (optimized) & $-0.02 \pm 0.01$ & --- & $0.051 \pm 0.008$ & --- \\
Propofol (GABAergic) & $-0.20 \pm 0.04$ & $-0.18 \pm 0.04$ & $0.029 \pm 0.006$ & $-0.022$ \\
Seizure (excitatory) & $+0.15 \pm 0.05$ & $+0.17 \pm 0.05$ & $0.021 \pm 0.007$ & $-0.030$ \\
5-MeO-DMT (serotonergic) & $+0.07 \pm 0.03$ & $+0.09 \pm 0.03$ & $0.068 \pm 0.011$ & $+0.017$ \\
\bottomrule
\end{tabular}
\end{table}

The modified van Albada--Robinson mean-field model \citep{vanalbada2009mean}, optimized to match empirical waking electrodynamics via Bayesian-genetic methods \citep{sherif2020bayesian}, demonstrated that proximity to edge-of-chaos criticality governs cross-frequency transfer efficiency (Table~\ref{tab:model}). The waking state was characterized by a Lyapunov exponent near zero ($\lambda = -0.02 \pm 0.01$), indicating dynamics just on the stable side of the critical point. Simulated propofol anesthesia (GABAergic modulation of synaptic decay rate $\alpha$) shifted dynamics away from the edge of chaos ($\Delta\lambda = -0.18$), reducing model cross-frequency TE by 0.022 bits. Simulated seizures pushed dynamics into the chaotic regime ($\lambda = +0.15$), also reducing TE. The psychedelic condition tuned dynamics closer to criticality from the chaotic side ($\Delta\lambda = +0.09$), enhancing TE by 0.017 bits.

\begin{table}[H]
\centering
\caption{Correlation between Lyapunov exponent proximity to zero and cross-frequency transfer entropy across 100 model parameter perturbations. Pearson correlations with 95\% bootstrap CIs.}
\label{tab:correlation}
\begin{tabular}{lccccc}
\toprule
Direction & $r$ & 95\% CI & $R^2$ & $F$ & $p$ \\
\midrule
Ctx $\to$ Thal & $-0.84$ & [$-0.89$, $-0.77$] & 0.71 & 238.4 & $< 0.001$ \\
Thal $\to$ Ctx & $-0.79$ & [$-0.85$, $-0.71$] & 0.62 & 162.1 & $< 0.001$ \\
\bottomrule
\end{tabular}
\end{table}

The absolute distance of the Lyapunov exponent from zero ($|\lambda|$) was strongly negatively correlated with cross-frequency transfer entropy across 100 parameter perturbations of the model (Ctx$\to$Thal: $r = -0.84$, $R^2 = 0.71$, $p < 0.001$; Thal$\to$Ctx: $r = -0.79$, $R^2 = 0.62$, $p < 0.001$; Table~\ref{tab:correlation}). This quantitative relationship formalizes the principle that proximity to edge-of-chaos criticality maximizes the capacity for cross-frequency information transfer.

\subsection{Summary of Statistical Comparisons}

\begin{table}[H]
\centering
\caption{Comprehensive summary of significance across all analyses and conditions. $^{*}p<0.05$, $^{**}p<0.01$, $^{***}p<0.001$. One-tailed Wilcoxon signed-rank tests for empirical data; Pearson correlation for model. ``Consistent'' indicates alignment with consciousness-level hypothesis.}
\label{tab:summary}
\begin{tabular}{lcccc}
\toprule
Analysis & Unconscious (Ctx$\to$T) & Unconscious (T$\to$Ctx) & Psychedelic (Ctx$\to$T) & Psychedelic (T$\to$Ctx) \\
\midrule
Cross-freq spectral TE & Decrease$^{**}$ & Decrease$^{*}$ & Increase$^{*}$ & NS \\
Broadband TE & NS & NS & NS & NS \\
Phase-amplitude coupling & NS & NS & NS & NS \\
Model ($|\lambda|$ vs.\ TE) & $r = -0.84^{***}$ & $r = -0.79^{***}$ & --- & --- \\
\bottomrule
\end{tabular}
\end{table}

\section{Discussion}

Our findings reveal a conserved spectral channel for bidirectional cortical--thalamic information transfer and identify edge-of-chaos criticality as the dynamical mechanism governing its modulation across states of consciousness.

\subsection{A Conserved Cross-Frequency Communication Channel}

The identification of a low-frequency-to-high-gamma transfer channel across four species, multiple thalamic nuclei, and distinct cortical areas establishes a fundamental organizing principle of thalamocortical communication. The slow oscillations (1.5--13~Hz) at the sender likely reflect the temporal structure of information packets, while high gamma activity (52--104~Hz) at the receiver may represent the local computational response to incoming information \citep{canolty2010importance}. This cross-frequency organization allows efficient multiplexing of information across temporal scales, enabling the integration of slowly varying contextual signals with rapidly fluctuating local computations \citep{aru2015untangling}.

\subsection{Spectral Resolution Is Essential}

The failure of broadband transfer entropy and phase-amplitude coupling to detect consciousness-related changes highlights a critical methodological point: spectrally resolved measures are necessary to capture the frequency-specific organization of neural communication. Broadband TE conflates contributions from multiple frequency bands, potentially canceling consciousness-related changes in specific channels with unrelated fluctuations in others. PAC, while frequency-resolved, lacks the directionality and information-theoretic grounding of spectral TE. These findings have important implications for the development of consciousness monitors and neural communication biomarkers.

\subsection{Edge-of-Chaos Criticality as a Unifying Mechanism}

The computational modeling results provide a mechanistic explanation for why consciousness states modulate cross-frequency transfer. The waking brain operates near the edge of chaos, where sensitivity, dynamic range, and information transmission are maximized \citep{shew2013functional}. Anesthesia pushes thalamocortical dynamics toward excessive stability (more negative $\lambda$), reducing the system's capacity to transmit information across frequencies. Seizures push dynamics into chaos (positive $\lambda$), where signal transmission becomes unreliable. Psychedelics tune dynamics closer to the critical point, potentially by modulating serotonergic gain, enhancing the capacity for cross-frequency communication \citep{carhart2016neural}.

This framework reconciles the apparently paradoxical spectral findings: both propofol and 5-MeO-DMT produce similar spectral power changes (increased slow power, decreased high-frequency power), yet they have opposing effects on directed information transfer because they shift dynamics in opposite directions relative to the critical point.

\subsection{Directional Asymmetry in Psychedelic Modulation}

The selective enhancement of cortex-to-thalamus transfer by 5-MeO-DMT, without significant enhancement of thalamus-to-cortex transfer, suggests that psychedelics may preferentially amplify top-down cortical influence on thalamic processing. This is consistent with models proposing that psychedelics increase the precision of top-down predictions \citep{carhart2016neural} and with the known distribution of serotonergic receptors in cortical output layers.

\subsection{Limitations}

Several limitations warrant consideration. First, the small sample sizes per species (5--8 subjects) limit statistical power, particularly for detecting the thalamus-to-cortex psychedelic effect. Second, the use of different thalamic nuclei and cortical areas across species, while strengthening the generalizability of the findings, introduces variability in the underlying connectivity. Third, the mean-field model, while biologically constrained and empirically optimized, necessarily simplifies the complexity of the basal ganglia--thalamocortical system. Fourth, the causal relationship between criticality and cross-frequency transfer, while supported by the model, remains correlational in the empirical data.

\section{Methods}

\subsection{Electrophysiology Recordings}

Simultaneous thalamic--cortical recordings were obtained from four preparations: (1) humans with essential tremor undergoing deep brain stimulation electrode implantation (Vim thalamus and sensorimotor cortex, $n = 8$), (2) Long-Evans rats (ventral posterior thalamus and somatosensory cortex, $n = 6$), (3) C57/BL6 mice (mediodorsal thalamus and medial prefrontal cortex, $n = 5$), and (4) GAERS rats (ventral posterior thalamus and contralateral somatosensory cortex, $n = 7$).

\subsection{Spectrally Resolved Transfer Entropy}

Directed information transfer was quantified using a surrogate-based method that measures the bits of transfer entropy lost when dynamics within specific frequency ranges are randomized. For each 10-second window, the time series was decomposed into frequency bands, and surrogates were generated by randomizing the phase of oscillations within the target frequency range while preserving all other spectral components. Transfer entropy was computed as the difference between the original and surrogate-randomized conditions.

\subsection{Brain State Manipulations}

Propofol anesthesia was administered to humans (target-controlled infusion to loss of consciousness) and Long-Evans rats (10~mg/kg bolus i.v.). Generalized spike-and-wave seizures occurred spontaneously in GAERS rats and were identified by characteristic 7--11~Hz spike-and-wave complexes. 5-MeO-DMT (10~mg/kg, i.p.) was administered to C57/BL6 mice.

\subsection{Computational Model}

A modified van Albada--Robinson mean-field model of the basal ganglia--thalamocortical system was implemented with sigmoidal firing rate functions $Q_a = Q_{\max} / (1 + \exp(-(V_a - \theta)/\sigma'))$ and synaptic response kernels parameterized by decay rate $\alpha$ and rise rate $\beta$. Model parameters were optimized via Bayesian-genetic methods to match empirical waking-state electrodynamics. Anesthesia was simulated by modulating the GABAergic synaptic decay rate $\alpha$ without altering maximal postsynaptic current. The Lyapunov exponent was computed from the linearized system dynamics to quantify proximity to edge-of-chaos criticality.

\subsection{Statistical Analysis}

Within-species, within-subject comparisons used one-tailed Wilcoxon signed-rank tests with effect sizes computed as $r = Z / \sqrt{n}$. Power spectral density was estimated using Welch's method \citep{welch1967use}. Model correlations used Pearson's $r$ with 95\% bootstrap confidence intervals (10,000 resamples). Significance threshold was $\alpha = 0.05$.

\section{Conclusion}

We demonstrate that a conserved cross-frequency spectral channel---low-frequency sender to high-gamma receiver---mediates bidirectional cortical--thalamic information transfer across species. This channel is selectively disrupted during unconscious states and enhanced by psychedelics. Computational modeling reveals that edge-of-chaos criticality of slow thalamocortical dynamics governs the efficiency of this communication, providing a unified mechanistic account of how brain dynamics shape consciousness through cross-frequency information transfer. These findings highlight the necessity of spectrally resolved, information-theoretic approaches for characterizing neural communication and suggest that the proximity to criticality may serve as a quantitative biomarker of conscious processing capacity.

\bibliographystyle{plainnat}
\bibliography{references}

\end{document}
